% !PRINTED 494
\weberpage{494}
%<<CONT
sind, daß darin aber im Falle $n = 7$ sämtliche $A_i$ rational von $f(\omega)^{24}$ abhängen, dagegen im Falle $n = 11$ die $A_i$ mit geradem Index ebenfalls rational von $f(\omega)^{24}$ abhängen, während die mit ungeradem Index das Produkt von $f(\omega)^{12}$ mit einer rationalen Funktion von $f(\omega)^{24}$ sind. Die Zahlkoeffizienten enthalten außer rationalen Zahlen nur noch die Irrationalität $i\sqrt{7}$ bzw. $i\sqrt{11}$, und demnach wollen wir diese Gleichungen bezeichnen durch

\weq{17}{\begin{aligned}
\Phi\bigl[w, f(\omega)^{24},\ \ i\sqrt{7}\bigr] &= 0,\\
\Phi\bigl[w, f(\omega)^{12},\ \ i\sqrt{11}\bigr] &= 0.
\end{aligned}}

Ändern wir gleichzeitig die Vorzeichen von $i$ und $w$, so ergiebt sich, wie wir in \S\,82 gesehen haben, eine Gleichung, deren Wurzeln die Größen $w'$ sind.

Auch bei der Anwendung der Substitution $(c''')$ zeigt sich für die beiden Fälle ein Unterschied.

Es ergiebt sich nämlich aus (13), mit Rücksicht darauf, daß nach (2) das Produkt sämtlicher $v_\nu$ den Wert $u^8$ hat, daß im Falle $n = 7$ durch die Substitution $(c''')$ die $w_\nu$ nur untereinander vertauscht werden, und zwar in folgender Weise:

\[
\begin{pmatrix} w_0, & w_1, & w_2, & w_3, & w_4, & w_5, & w_6\\ w_5, & w_3, & w_6, & w_2, & w_4, & w_0, & w_1 \end{pmatrix},
\]

und daraus folgt, daß in diesem Falle die Koeffizienten $A_\nu$ alle ungeändert bleiben, durch die Vertauschung:

\[
\left(f(\omega)^{24},\ \frac{2^{12}}{f(\omega)^{24}}\right).
\]

Überdies kommen im Nenner dieser Koeffizienten nur Potenzen von $f(\omega)^{24}$ vor (\S\,73), und die Potenzen von $f(\omega)^{24}$ steigen bis zu derselben Höhe wie die von $f(\omega)^{-24}$.

Im Falle $n = 11$ gehen durch die Vertauschung $(c''')$ die Größen $w_\nu$ in die Größen $w'_\nu$ über, und zwar in folgender Reihenfolge:

\[
\begin{pmatrix} w_0, & w_1, & w_2, & w_3, & w_4, & w_5, & w_6, & w_7, & w_8, & w_9, & w_{10}\\ w'_8, & w'_2, & w'_7, & w'_0, & w'_9, & w'_5, & w'_6, & w'_4, & w'_3, & w'_{10}, & w'_1 \end{pmatrix}.
\]

Nach (17) können wir diese Eigenschaft durch die in bezug auf $w$ identische Gleichung ausdrücken:

\[
-\Phi\bigl[-w, f(\omega)^{12}, -i\sqrt{11}\bigr] = \Phi\!\left(w, \frac{2^6}{f(\omega)^{12}}, i\sqrt{11}\right),
\]

und daraus ergeben sich die folgenden Eigenschaften der Koeffizienten $A_\nu$:
%CONT>>
